% Shared layout and graphics settings.
\graphicspath{{assets/images/}{cover/}}

\tolerance        = 1000
\emergencystretch = 0.8em

\widowpenalty        = 10000
\clubpenalty         = 10000
\displaywidowpenalty = 10000

\XeTeXlinebreaklocale "zh"
\XeTeXlinebreakskip = 0pt plus 0.2pt minus 0.1pt

% --- 行距：中文五号正文的舒适区间为字号的 1.5–1.6 倍 ---
% 10pt 基础下 \normalsize 基线 12pt × 1.35 = 16.2pt；
% 配合 CJK Scale=1.05（正文五号 10.5pt）→ 16.2/10.5 ≈ 1.54×。
% 思源宋体字面率大于霞鹜新致宋，自初版 1.3 上调以维持视觉行距。
% \linespread 需在 \begin{document} 处 \selectfont 生效，以防被类/ctex 覆盖。
\linespread{1.35}
\AtBeginDocument{%
  \selectfont
  % 使用当前中文字体的字宽，避免 1.05 倍 CJK 缩放造成缩进不足两字。
  \setlength{\parindent}{2\ccwd}%
  \raggedbottom
}

\setlength{\parskip}{0pt plus 0.5pt}
% 每次字号命令完成后重新设定公式间距；使用 etoolbox 补丁兼容旧版类文件。
\newcommand{\SeriesDisplaySpacing}{%
  \setlength{\abovedisplayskip}{8pt plus 2pt minus 2pt}%
  \setlength{\belowdisplayskip}{8pt plus 2pt minus 2pt}%
  \setlength{\abovedisplayshortskip}{3pt plus 2pt}%
  \setlength{\belowdisplayshortskip}{6pt plus 2pt minus 2pt}%
}
\apptocmd{\normalsize}{\SeriesDisplaySpacing}{}{}
\apptocmd{\small}{\SeriesDisplaySpacing}{}{}
\apptocmd{\footnotesize}{\SeriesDisplaySpacing}{}{}
\apptocmd{\scriptsize}{\SeriesDisplaySpacing}{}{}
\apptocmd{\tiny}{\SeriesDisplaySpacing}{}{}
\apptocmd{\large}{\SeriesDisplaySpacing}{}{}
\apptocmd{\Large}{\SeriesDisplaySpacing}{}{}
\apptocmd{\LARGE}{\SeriesDisplaySpacing}{}{}
\apptocmd{\huge}{\SeriesDisplaySpacing}{}{}
\apptocmd{\Huge}{\SeriesDisplaySpacing}{}{}
\AtBeginDocument{\SeriesDisplaySpacing}

% 覆盖类的 nolistsep；保留各层编号及正文局部选项。
\setlist{
  topsep=6pt plus 2pt minus 2pt,
  itemsep=3pt plus 1pt minus 1pt,
  parsep=0pt,
  partopsep=0pt
}
\setlist[enumerate,1]{leftmargin=2.4em, labelsep=0.5em}
\setlist[itemize,1]{leftmargin=2em, labelsep=0.5em}

% 图表与文字留出半行以上的距离，避免强制齐底拉大段间空白。
\setlength{\textfloatsep}{12pt plus 3pt minus 2pt}
\setlength{\intextsep}{10pt plus 2pt minus 2pt}
\setlength{\floatsep}{10pt plus 2pt minus 2pt}

\renewcommand*\arraystretch{1.3}
